\section{Canonical High-Correlation Comparisons}
\label{app:verification}\label{app:tail}

The middle-range argument ends at \(\ell=7/2\). The following
three comparisons cover the entire remaining half-line with the same
canonical penalty \(\gamma=\lambda\alpha\). The spectral multiplier
is \(p=(\lambda+1/\alpha)^2\); it is not replaced by the angle
constant \(P=\pi^2/4\). The constant \(P\) still denotes the
limit of the angle profile \(c(v)\); it supplies elementary upper
bounds and measures the angle saving at the half-height join. Keeping
the exact spectral multiplier \(p\) preserves the low-degree
cancellations in the join and cutoff proofs.

\subsection{Reduction to Three Comparisons}
\label{app:scalar-map}

Write \(r=\rho=\tanh\ell\), \(x=e^{-\ell}\), \(L=\log2\), and
\(\theta=\arcsin r\). Our definitions are
\begin{equation}
 \begin{split}
 \lambda&=\theta/r,qquad
 \alpha=(\theta+\sinh\ell)/\ell,\\
 \gamma&=\lambda\alpha,qquad d=\ell/2+5/4.
 \end{split}
 \label{eq:common-policy}
\end{equation}
As before, \(\weight(s)=\gamma s/(\gamma+s)\),
\(d_1=[\weight(d)-\weight(1)]r^2\),
\(d_2=[\weight(d)-\weight(2)]r^4\), and \(d_b=d_1-d_2/2\).
The retained-variance baseline is
\begin{align*}
 p&=(\lambda+1/\alpha)^2,& b&=\lambda^2(2+1/\gamma),\\
 C_d&=p\weight(d)-b,& B&=2\theta^2+C_d\phi(r)/L.
\end{align*}
Here \(C_d\) is a channel-dependent spectral coefficient, not an
entropy approximation constant. The degree balances the LSI exponential
against the entropy ratio at half the channel height; the penalty
itself is already fixed by the dictator-centered scalar tangent.

The local lemma handles influences at least
\(a_{\rm cut}=1/[1+(8/3)x]\). For the remaining coordinates we prove:
\begin{enumerate}
\item The line-angle LSI yield is below \(B/\ell\) for all heights
up to \(z=\ell/2\) (Section~\ref{app:manual-small}).
\item The original Parseval yield is below \(B/\ell\) at the same
join (Section~\ref{app:manual-join-direct}).
\item That yield is below \(B/\ell\) at the cutoff height
\(F(V)=F(\ell)/a_{\rm cut}\)
(Section~\ref{app:manual-cutoff-common}).
\end{enumerate}
We first check the lower-variance sign, all-bias mean condition, and
coefficient needed for the global endpoint principle. That principle
then fills the interval from \(z\) to \(V\); if \(V\le z\), the
small-height comparison already covers it. No channel subdivision is
needed in any of these checks.

\subsection{Local Cutoff}
\label{app:manual-local}\label{app:analytic-local}

Write \(x=e^{-\ell}\). We use the single influence cutoff
\begin{equation}
                 a_{\rm cut}=\frac{1}{1+(8/3)x}.
 \label{eq:rational-local-cutoff}
\end{equation}
We need this cutoff only for \(\ell\ge7/2\). On this half-line
\(x<1/32\): indeed \(e^3>20\) and
\(e^{1/2}>1+1/2+1/8=13/8\) give \(e^{7/2}>32\).
The proof retains the leading entropy logarithms until they cancel,
then bounds the three remaining errors uniformly.

Put \(r=(1-x^2)/(1+x^2)\), and let
\(h_b(p)=-p\log p-(1-p)\log(1-p)\) denote binary entropy in
probability coordinates. The following temporary quantities expose
the cancellation between the two entropies:
\begin{align*}
 \kappa_{\rm loc}&=\frac4{3+8x},&
 \eta_{\rm loc}&=\frac{3x}{4(1+x^2)},\\
 p&=\frac{x^2}{1+x^2},&
 q_{\rm loc}&=\kappa_{\rm loc}x(1+\eta_{\rm loc}),\\
 A_{\rm loc}&=\frac1{2x^2}+\frac3{4x}+1+\frac{x^2}{2}.
\end{align*}
Thus \(p=(1-r)/2\), \(q_{\rm loc}=(1-r a_{\rm cut})/2\). The local margin
\((1-r^2)h(r a_{\rm cut})-(1-a_{\rm cut}r^2)h(r)\), divided by the positive factor
\(4\kappa_{\rm loc} x^3/(1+x^2)^2\), equals
\[
 \Lambda_{\rm loc}=\frac{h_b(q_{\rm loc})}{\kappa_{\rm loc}x}
 -A_{\rm loc}h_b(p).
\]

\paragraph{Cancellation of the leading logarithms.}
Both entropy estimates follow from one pair of logarithm bounds:
\[
 \frac{2t}{2+t}\le\log(1+t)\le t-\frac{t^2}{2(1+t)}
 \qquad(t\ge0).
\]
These are the midpoint and trapezoid bounds for the integral of the
convex function \(s\mapsto(1+ts)^{-1}\) on \([0,1]\), multiplied by \(t\).
The lower bound, with \(t=q_{\rm loc}/(1-q_{\rm loc})\), gives
\begin{align*}
 h_b(q_{\rm loc})\ge{}&q_{\rm loc}\log(1/q_{\rm loc})+q_{\rm loc}\\
 &-\frac{q_{\rm loc}^2}{2-q_{\rm loc}}.
\end{align*}
The upper bound gives
\((1+\eta_{\rm loc})\log(1+\eta_{\rm loc})\le\eta_{\rm loc}+\eta_{\rm loc}^2/2\).
Consequently
\begin{align*}
 \frac{h_b(q_{\rm loc})}{\kappa_{\rm loc}x}\ge{}&
 (1+\eta_{\rm loc})(\ell-\log\kappa_{\rm loc})+1\\
 &-\frac{\eta_{\rm loc}^2}{2}
 -\frac{\kappa_{\rm loc}x(1+\eta_{\rm loc})^2}{2-q_{\rm loc}}.
\end{align*}
Here \(\eta_{\rm loc}\le3x/4\) and
\(q_{\rm loc}\le x(4+3x)/(3+8x)\). Keep their two errors together:
\begin{align*}
 &\frac{\eta_{\rm loc}^2}{2}
 +\frac{\kappa_{\rm loc}x(1+\eta_{\rm loc})^2}{2-q_{\rm loc}}\\
 &\quad\le\frac9{32}x^2+\frac{x(4+3x)^2}{12(2+4x-x^2)}
 \le\frac23x+\frac15x^2.
\end{align*}
For the last inequality, the identity
\[
 (8-x)(2+4x-x^2)-(4+3x)^2=x(6-21x+x^2)>0
\]
holds for \(x<1/4\). It bounds the rational term by
\(2x/3-x^2/12\); then \(9/32-1/12<1/5\).

For the other entropy, write
\(h_b(p)=\log(1+x^2)+2\ell x^2/(1+x^2)\).
The same trapezoid bound for the logarithm yields
\begin{align*}
 A_{\rm loc}\log(1+x^2)
 &\le\frac12+\frac34x+\frac34x^2\\
 &\quad+\frac{x^3[2x(1+x^2)-3]}{8(1+x^2)}\\
 &\le\frac12+\frac34x+\frac34x^2.
\end{align*}
The remainder is negative for \(x<1/4\). Therefore
\begin{align*}
 A_{\rm loc}h_b(p)\le{}&\frac12+\frac34x+\frac34x^2\\
 &+\frac{2A_{\rm loc}\ell x^2}{1+x^2}.
\end{align*}
The coefficient of \(\ell\) now cancels exactly to
\begin{align*}
 1+\eta_{\rm loc}-\frac{2A_{\rm loc}x^2}{1+x^2}
 &=-x\left[\frac{3}{4(1+x^2)}+x\right].
\end{align*}
Therefore
\begin{equation}
 \begin{split}
 \Lambda_{\rm loc}\ge{}&\frac12-(1+\eta_{\rm loc})\log\kappa_{\rm loc}
 -\frac{17}{12}x-\frac{19}{20}x^2\\
 &-\ell x\left[\frac{3}{4(1+x^2)}+x\right].
 \end{split}
 \label{eq:manual-local-cancellation}
\end{equation}

\paragraph{Uniform remainder estimates.}
The same trapezoid bound at \(t=1/3\) gives
\(\log(4/3)<7/24\). Since \(\eta_{\rm loc}\le3/128\)
and \(\kappa_{\rm loc}<4/3\),
\[
 (1+\eta_{\rm loc})\log\kappa_{\rm loc}
 <\left(1+\frac3{128}\right)\frac7{24}<\frac13.
\]
Also
\[
 \frac{17}{12}x+\frac{19}{20}x^2
 <\frac{17}{12\cdot32}+\frac{19}{20\cdot32^2}<\frac1{20}.
\]
Finally \(\ell e^{-\ell}\) decreases here, so \(\ell x<7/64\), and
\[
 \ell x\left[\frac3{4(1+x^2)}+x\right]
 <\frac7{64}\left(\frac34+\frac1{32}\right)<\frac3{32}.
\]
Substitution in \eqref{eq:manual-local-cancellation} proves
\begin{align*}
 \Lambda_{\rm loc}&>\frac12-\frac13-\frac1{20}-\frac3{32}\\
 &=\frac{11}{480}>0.
\end{align*}
Thus the criterion holds at \(a_{\rm cut}\) throughout
\(\ell\ge7/2\). For fixed \(r\), its left side
\((1-r^2)h(ra)-(1-ar^2)h(r)\) is concave in \(a\) and vanishes
at \(a=1\). It is therefore nonnegative for every
\(a\in[a_{\rm cut},1]\).

\subsection{Canonical Penalty and Preliminary Bounds}
\label{app:canonical-prerequisites}

Throughout this appendix \(\ell\ge7/2\), \(x=e^{-\ell}\), and
\(\theta=\arcsin r\). The rule is
\begin{align*}
 \lambda&=\theta/r,& \alpha&=(\theta+\sinh\ell)/\ell,\\
 \gamma&=\lambda\alpha,& d&=\ell/2+5/4,
\end{align*}
and
\begin{align*}
 p&=(\lambda+1/\alpha)^2,& b&=\lambda^2(2+1/\gamma),\\
 C_d&=p\weight(d)-b,& B&=2\theta^2+C_d\phi(r)/L.
\end{align*}
The elementary bounds \(e^{7/2}>33\), \(157/50<\pi<22/7\), and
\(2/3<L<3/4\) give \(r>99/100\), \(\theta>3/2\), and
\(\lambda^2<P<5/2\); the angle ratio is bounded by the endpoint
chord of the convex function \(\arcsin r\). Since
\(\theta+\sinh\ell>1/(2x)\),
\begin{equation}
 \gamma>\frac{3e^\ell}{4\ell}>\ell+\frac72=2d+1.
 \label{eq:canonical-penalty-lower}
\end{equation}
For the second inequality, \(3e^\ell-4\ell^2-14\ell\) is positive
at \(7/2\), its first derivative is positive there, and its second
derivative \(3e^\ell-8\) stays positive.

The exact multiplier identities are
\begin{equation}
 \begin{split}
 p\weight(1)-b&=-\lambda^2,\\
 p\weight(2)-b&=-\frac{\lambda^2}{\gamma+2},\\
 C_d&=\lambda^2\left[d-2-\frac{(d-1)^2}{\gamma+d}\right].
 \end{split}
 \label{eq:canonical-multipliers}
\end{equation}
Thus \(C_d>0\), since \(\gamma(d-2)>1\), and
\(C_d<(5/2)(d-2)\), \(b<(5/2)(2+1/7)<6\).
The all-bias mean condition follows without a separate channel test:
\begin{align*}
 \gamma-b-C_d\left(1-\frac1{2L}\right)
 &>\ell+\frac72-6-\frac56\left(\frac\ell2-\frac34\right)\\
 &=\frac16+\frac7{12}\left(\ell-\frac72\right)>0.
\end{align*}
Here \(0<1-1/(2L)<1/3\). Finally, the coefficient required by the
original-yield endpoint principle satisfies
\begin{align*}
 \frac{p d_2}{2r}
 &=\frac p2(d-2)\frac\gamma{\gamma+d}
                   \frac\gamma{\gamma+2}r^3\\
 &>1\cdot1\cdot\frac12\cdot\frac12\cdot\frac34
 =\frac3{16}>\frac5{32}.
\end{align*}
Also \(0<d_2<d_1\), so \(d_b>0\). The global shape inequality
\eqref{eq:original-shape} therefore applies with multiplier \(p\).

\paragraph{Entropy estimate.}
Put \(E=h(r)-L(1-r^2)\ge0\). For every \(\ell\ge3/2\),
\begin{equation}
                     E\le x^2(2\ell-5/3).
 \label{eq:shared-entropy-error}
\end{equation}
Indeed, with \(q=x^2\), the entropy formula gives
\begin{align*}
 2\ell-\frac E q
 &=\frac{2\ell q}{1+q}+\frac{4L}{(1+q)^2}
 -\frac{\log(1+q)}q.
\end{align*}
The reciprocal trapezoid bound gives
\(\log(1+q)/q\le(2+q)/[2(1+q)]\). Using \(2\ell\ge3\) leaves
\[
 2\ell-\frac E q-\frac53
 \ge\frac{4(L-2/3)-11q/6+5q^2/6}{(1+q)^2}>0.
\]
The last numerator exceeds \(8/81-11/120>0\): the first two terms
of \(L=2\operatorname{atanh}(1/3)\) give \(L-2/3>2/81\), and
\(e^3>20\) gives \(q<1/20\).

\subsection{Small Heights: Paired Spectral Costs}
\label{app:manual-small}

Let \(z=\ell/2\). The line-angle LSI comparison is
\begin{align*}
 M_{\rm small}&=\frac B\ell-r
       -\frac{pK_{\rm LS}r^2F(\ell)}{zF(z)},\\
 K_{\rm LS}&=\frac{\gamma^2 e^{2d-3}}{2(\gamma+d)(\gamma+1)}.
\end{align*}
Set \(a_d=d-1\), \(k=e^{-1/2}\),
\(T=e^\ell F(\ell)/F(z)\), and \(f_d=C_d/\lambda^2\).
The identities \eqref{eq:canonical-multipliers} and
\(\phi(r)/L=r^2-E/L\) give the useful cancellation
\begin{equation}
 \begin{split}
 \frac{\ell M_{\rm small}}{\theta^2}
 ={}&d-\frac\ell{c(\ell)}-kT\\
 &-\frac{a_d(a_d-kT)}{\gamma+d}-\frac{f_dE}{Lr^2}.
 \end{split}
 \label{eq:canonical-small-pairing}
\end{equation}
We bound the first three terms below, keeping the resolvent loss
paired with the profile cost in the fourth term.

The elementary logarithm bounds
\(t/(1+t)\le\log(1+t)\) and
\((1+t)\log(1+t)\le t+t^2/2\), obtained by integration, yield
\begin{align*}
 \frac{2\ell+1}{(\ell+1+x/2)(1+x)}
 &\le T\le\frac{2\ell+1+x^2/2}{(\ell+1)(1+x)}\\
 &\le2-\frac1{\ell+1}.
\end{align*}
The lower bound exceeds \(5/3\): after multiplication by its
positive denominator, the difference has numerator
\[
 \ell-2-5\ell x-\frac{15}2x-\frac52x^2
 \ge\frac32-\frac{25}{32}-\frac5{2048}>0.
\]
Here this expression increases with \(\ell\) and decreases with
\(x\) on \(\ell\ge7/2\), \(0<x<1/32\).
Since \(3/5<k<5/8\) and \(c(\ell)>\theta^2>9/4\),
\begin{align*}
 1&<kT<\frac58\left(2-\frac1{\ell+1}\right),\\
 d-\frac\ell{c(\ell)}-kT
 &>\frac\ell{18}+\frac5{8(\ell+1)}\ge\frac13.
\end{align*}
The last difference from \(1/3\) equals
\((2\ell-3)(2\ell-7)/[72(\ell+1)]\).

Next, \(\gamma+d>4(d-1)(d-2)\). By
\eqref{eq:canonical-penalty-lower}, it suffices that
\(3e^\ell>4\ell^3-6\ell^2-8\ell\).
For \(s=\ell-7/2\ge0\), the positive cubic Taylor sum for \(e^s\)
and \(e^{7/2}>33\) bound this difference below by
\[
                  29+2s+\frac{27}2s^2+\frac{25}2s^3>0.
\]
Consequently
\[
 \frac{a_d(a_d-kT)}{\gamma+d}
 <\frac{a_d(a_d-1)}{\gamma+d}<\frac14.
\]
Finally, \eqref{eq:shared-entropy-error}, \(f_d<d-2\),
\(L>2/3\), and \(r^2>3/4\) give
\[
 \frac{f_dE}{Lr^2}<2(d-2)(2\ell-5/3)e^{-2\ell}<\frac1{96}.
\]
The product decreases: its logarithmic derivative is at most
\(1/2+3/8-2<0\). At \(7/2\) its value is less than
\(2(16/3)/32^2=1/96\). Thus
\[
 M_{\rm small}>\frac{\theta^2}{\ell}
                      \left(\frac13-\frac14-\frac1{96}\right)
                =\frac{7\theta^2}{96\ell}>0.
\]
Since \(vF(v)\) decreases and \(c(v)\le v\), this comparison at
\(z\) covers every \(0<v\le z\). In particular \(B/\ell>r\).

\subsection{Half-Height Join and Angular Saving}
\label{app:joining-height}\label{app:manual-join-direct}

Write \(a_z=F(\ell)/F(z)\). The original Parseval margin at the
join has the exact form
\begin{align*}
 M_{\rm join}
 &:={Bz}/{\ell}-r c(z)-p[d_2/2+d_ba_z]\\
 &=r[c(\ell)-c(z)]-\frac{\lambda^2r^4}{2(\gamma+2)}
   +\frac{C_d}{2}\left[1-r^4-\frac{h(r)}L\right]\\
 &\quad-p d_ba_z.
\end{align*}
The multiplier identities also give
\begin{equation}
 \begin{split}
 p d_b={}&C_d(r^2-r^4/2)
       +\lambda^2\left[r^2-\frac{r^4}{2(\gamma+2)}\right]\\
       &\le C_d/2+P\le Pd/2.
 \end{split}
 \label{eq:canonical-db-pairing}
\end{equation}

The angle saving is of order \(\sqrt x\), while the losses are of
order \(x\). To see it directly, the function
\(e^v[P-c(v)]\) increases. With \(t=\arcsin(\tanh v)\),
\begin{align*}
 &(t+\cos t)^2-c(v)-c'(v)\\
 &\quad=(1-\sin t)(t/\sin t-\cos t)^2
 +\sin t\cos^2t>0.
\end{align*}
As \(t+\cos t<\pi/2\), this gives \(c+c'<P\).
At \(v=\log(10/3)\), the bounds \(4/\pi>5/4\) and
\(\arctan y\ge y-y^3/3\) at \(y=3/10\) give
\[
 2t/\pi<16/25,\qquad
 c(v)/P<(16/25)^2(109/91)<1/2.
\]
Hence \(P-c(z)>(5P/3)\sqrt x\), since
\(z\ge7/4>\log(10/3)\). On the other hand,
\(c(\ell)\ge\theta^2\ge P-2\pi x\), so
\[
                 c(\ell)-c(z)>\frac{5P}{3}\sqrt x-2\pi x.
\]
Also \(h(r)<(2\ell+1)x^2\), and
\(\log(1+x)\ge x/(1+x)\) gives
\[
 \frac{F(\ell)}{F(z)}
 <x\frac{2\ell+1}{\ell+1}\frac{1+x^2}{1+x}<2x.
\]
Using these bounds, \eqref{eq:canonical-db-pairing}, and
\eqref{eq:canonical-penalty-lower}, and discarding the positive
\(C_d(1-r^4)/2\), gives
\begin{align*}
 \frac{M_{\rm join}}{P\sqrt x}
 >{}&\frac{5r}{3}
 -\left(\frac76\ell+\frac54+\frac8\pi\right)e^{-\ell/2}\\
 &-\frac{(d-2)(2\ell+1)}{2L}e^{-3\ell/2}.
\end{align*}
The first term increases and both losses decrease. For the affine
loss this follows by differentiation; for the last, the logarithmic
derivative is
\(1/(\ell-3/2)+1/(\ell+1/2)-3/2<0\).
At \(\ell=7/2\), use \(r>99/100\), \(\sqrt x<7/40\),
\(8/\pi<8/3\), and \(2L>4/3\). The affine coefficient is less
than eight, and the last loss is less than \(6(7/40)^3<1/25\).
Thus the whole half-line satisfies
\[
             \frac{M_{\rm join}}{P\sqrt x}
                  >\frac{33}{20}-\frac75-\frac1{25}
                   =\frac{21}{100}>0.
\]

\subsection{Cutoff Bound with Exact Low-Degree Multipliers}
\label{app:manual-cutoff-common}\label{app:pivotal-tail-extension}

The local cutoff is \(a_{\rm cut}=1/[1+(8/3)x]\). Let
\(F(V)=F(\ell)/a_{\rm cut}\). The global profile bound
\(-F'/F>8/5\) and \(\log(1+t)<t\) give
\[
                  0<\ell-V<\frac53x.
\]
We check the actual Parseval margin
\[
 M_{\rm cut}=\frac B\ell V-r c(V)
                  -p[d_2/2+d_ba_{\rm cut}].
\]
By \eqref{eq:canonical-multipliers}, its exact dictator deficit is
\[
 B-pd_1-r c(\ell)=C_d[\phi(r)/L-r^2]=-C_dE/L.
\]
Concavity of \(c\), together with \(c'>0\), bounds the displacement
loss below by \(-(B/\ell)(\ell-V)\). Since \(B>0\),
\[
 M_{\rm cut}\ge-C_dE/L+p d_b(1-a_{\rm cut})
                              -\frac B\ell(\ell-V).
\]
This avoids any separate sign condition on a tangent coefficient.
Use \(B\le2\lambda^2r^2+C_d\),
\eqref{eq:canonical-db-pairing} as an equality in its first line,
and \eqref{eq:shared-entropy-error}. They give
\begin{align*}
 \frac{M_{\rm cut}}x>{}&C_d\left[
 \frac{8r^2(1-r^2/2)}{3+8x}-\frac{(2\ell-5/3)x}{L}
                                      -\frac5{3\ell}\right]\\
 &+\lambda^2r^2\left[
 \frac8{3+8x}\left(1-\frac{r^2}{2(\gamma+2)}\right)
                                      -\frac{10}{3\ell}\right].
\end{align*}
Both brackets have uniform positive lower bounds.
Indeed \(x<1/32\) implies \(r^2>3/4\), so
\[
 \frac{8r^2(1-r^2/2)}{3+8x}>\frac{32}{13}\frac{15}{32}>1.
\]
The function \((2\ell-5/3)e^{-\ell}\) decreases; using \(L>2/3\),
\[
 \frac{(2\ell-5/3)x}{L}<\frac{16}{3}\frac1{32}\frac32=\frac14,
 \qquad \frac5{3\ell}\le\frac{10}{21}<\frac12.
\]
Thus the first bracket exceeds \(1/4\). For the second,
\(8/(3+8x)>2\), \(1-r^2/[2(\gamma+2)]>3/4\), and
\(10/(3\ell)\le20/21<1\). It therefore exceeds \(1/2\).
We conclude on the entire half-line that
\[
              M_{\rm cut}>x\left(\frac{C_d}{4}
                                  +\frac{\lambda^2r^2}{2}\right)>0.
\]
This completes the third comparison. The original-yield endpoint
principle now fills the retained heights between the join and the
cutoff; heights below the join were covered by the LSI comparison.



\subsection{Fixed Constants}
\label{app:fixed-arithmetic}

All fixed exponential comparisons follow from one positive sum and
its geometric remainder: for \(0<y<14\),
\[
 E_{12}(y):=\sum_{j=0}^{12}\frac{y^j}{j!}
 <e^y<E_{12}(y)+\frac{y^{13}}{13!(1-y/14)}.
\]
This proves the exponential brackets used above by rational substitution.

For logarithms use, for \(|y|<1\),
\begin{align*}
 \log\frac{1+y}{1-y}
 &=2\sum_{j=0}^{N-1}\frac{y^{2j+1}}{2j+1}+R_N(y),\\
 |R_N(y)|&\le\frac{2|y|^{2N+1}}{(2N+1)(1-y^2)}.
\end{align*}
Taking \(y=1/3\) bounds \(\log2\); taking
\(y=(q-1)/(q+1)\) bounds \(\log q\) for every positive rational
\(q\). The low-correlation arithmetic uses \(N=12\).
The basic constant bounds are
\[
                 3.14159<\pi<3.1416,\qquad .69314<\log2<.69315.
\]
For \(\pi\), use Machin's identity
\[
 \pi=16\arctan(1/5)-4\arctan(1/239).
\]
Bound the first arctangent
between its first four and first five alternating terms, and the
second between its first two and first three. To verify the identity,
the tangent addition formula gives
\[
 \tan(4\arctan(1/5))=120/119;
\]
subtracting \(\arctan(1/239)\)
gives tangent one and an angle in \((0,\pi/2)\).
For \(\log2\), take
the first five terms of the displayed logarithm series at \(y=1/3\)
and bound the remainder by
\[
 \frac{2(1/3)^{11}}{11(1-1/9)}.
\]
In particular these imply the bounds used here,
\(25/8<\pi<22/7\), \(9/4<P<5/2\), and
\(2/3<L<3/4\). The positive exponential sums give
\(e^3>20\), \(e^{7/2}>33\), and
\(3/5<e^{-1/2}<5/8\).
Square-root bounds follow by squaring positive rational endpoints.

\paragraph{Completion.}
The three comparisons, the mean condition, and the local cutoff cover
every \(\ell\ge7/2\), including the shared endpoint. All estimates
are printed above; no interval receipt, symbolic sign certificate, or
external computation supplies a missing comparison.
